\scanpage{88}

which can be made explicit in terms of (78) by noting that every flag of $J\cup L$ \inserted{occurring in (78)} can be written uniquely
\[
(80)\qquad(j_{\unsure{1|q}}\mathord{\cdot}\mathord{-}\mathord{\cdot}j_p\ l_1\mathord{\cdot}\mathord{-}l_q)\qquad
\begin{gathered}[t]
i_{\unsure{q|0}}\mathord{-}\mathord{-}i_p\in J\\
l_1\mathord{-}\mathord{-}l_q\in L\\
p,q\geq1
\end{gathered}
\]
and \deleted{\unsure{then}} we have
\diagram{\begin{tikzcd}[ampersand replacement=\&,column sep=large,row sep=large]
(81)\&\Sigma_{j_1\mathord{-}\mathord{-}j_p\,l_1\mathord{-}\mathord{-}l_q}\arrow[r,hook]\arrow[d]\&\Sigma_{l_1\mathord{\cdot}\mathord{\cdot}\mathord{-}l_q}\\
\&\Sigma_{j_1,\mathord{\cdot}\mathord{-},j_p}\&
\end{tikzcd}}

\uline{Remark} The $\Sigma_{J,L}$ make sense \inserted{(in a suitable isotopy category)} independently of the local conical retractions around the $X_i^*$, at least in tame topological theory, as do the $\Sigma_{J,L}\hookrightarrow\Sigma_L$ --- and the description I have given of them here is adapted to this case --- the structure $\Sigma_{J,L}\to\Sigma_J$ on the other hand uses the local \deleted{conical} \deleted{structures} \inserted{fibrations} \inserted{\deleted{\unsure{of ri.}}}. When one has $=$ local conical structures (which, in tame topology, and in the equisingular case, is an additional ``harmless'' datum), it is more convenient,
\margin{87}
